\documentclass[aspectratio=169,11pt]{beamer}

\usepackage[T1]{fontenc}
\usepackage[utf8]{inputenc}
\usepackage[english]{babel}
\usepackage{graphicx}
\usepackage{booktabs}
\usepackage{array}
\usepackage{amsmath, amssymb}
\usepackage{xcolor}
\usepackage{hyperref}

\definecolor{usdred}{RGB}{161,7,43}
\definecolor{darkgray}{RGB}{45,45,45}

\usetheme{Madrid}
\usecolortheme{default}

\setbeamercolor{palette primary}{bg=usdred,fg=white}
\setbeamercolor{palette secondary}{bg=usdred,fg=white}
\setbeamercolor{palette tertiary}{bg=darkgray,fg=white}
\setbeamercolor{structure}{fg=usdred}
\setbeamercolor{frametitle}{bg=usdred,fg=white}
\setbeamercolor{title}{fg=usdred}

\setbeamertemplate{navigation symbols}{}
\setbeamertemplate{caption}[numbered]

\graphicspath{{outputs/figures/}{./}}

\title[Room Occupancy Detection]{Room Occupancy Detection Using Environmental Sensor Data}
\author[Benjamin van Rhee]{Benjamin van Rhee}
\institute[USD]{Department of Computer Science\\University of South Dakota\\CSC 444}
\date{Spring 2026}

\begin{document}

\begin{frame}
    \titlepage
\end{frame}

\begin{frame}{Project Overview}
    \begin{itemize}
        \item Goal: predict whether a room is occupied using environmental sensor data.
        \item Dataset: UCI Occupancy Detection dataset.
        \item Task type: supervised binary classification.
        \item Models compared:
        \begin{itemize}
            \item Logistic Regression
            \item Random Forest
        \end{itemize}
        \item Main result: scaled Logistic Regression performed best overall.
    \end{itemize}
\end{frame}

\begin{frame}{Why Occupancy Detection Matters}
    \begin{itemize}
        \item Smart buildings need reliable occupancy estimates.
        \item Occupancy information can improve:
        \begin{itemize}
            \item lighting control
            \item heating and cooling
            \item ventilation
            \item energy efficiency
        \end{itemize}
        \item Environmental sensors are less invasive than cameras.
        \item This makes the problem useful for offices, schools, and homes.
    \end{itemize}
\end{frame}

\begin{frame}{Dataset}
    \begin{itemize}
        \item Source: UCI Machine Learning Repository.
        \item Total observations: 20,560.
        \item Target variable: \texttt{Occupancy}.
        \item Sensor features:
        \begin{itemize}
            \item Temperature
            \item Humidity
            \item Light
            \item CO$_2$
            \item HumidityRatio
        \end{itemize}
        \item The timestamp column was parsed but not used as a model feature.
    \end{itemize}
\end{frame}

\begin{frame}{Dataset Splits}
    \begin{table}
        \centering
        \caption{Predefined Dataset Splits}
        \begin{tabular}{lr}
            \toprule
            Split & Rows \\
            \midrule
            Training & 8,143 \\
            Validation & 2,665 \\
            Final holdout test & 9,752 \\
            \bottomrule
        \end{tabular}
    \end{table}

    \vspace{0.3cm}
    The validation split was used for model comparison. The final holdout test split was used only for final evaluation.
\end{frame}

\begin{frame}{Class Balance}
    \begin{figure}
        \centering
        \includegraphics[width=0.78\textwidth]{class_balance.png}
        \caption{Class balance across training, validation, and test splits.}
    \end{figure}
\end{frame}

\begin{frame}{Class Balance Counts}
    \begin{table}
        \centering
        \caption{Occupancy Counts by Split}
        \begin{tabular}{lrr}
            \toprule
            Split & Not Occupied & Occupied \\
            \midrule
            Training & 6,414 & 1,729 \\
            Validation & 1,693 & 972 \\
            Test & 7,703 & 2,049 \\
            \bottomrule
        \end{tabular}
    \end{table}

    \vspace{0.3cm}
    Because the dataset is imbalanced, F1 score, precision, and recall are more informative than accuracy alone.
\end{frame}

\begin{frame}{Preprocessing}
    \begin{itemize}
        \item The first column in each data file was treated as a row index.
        \item No missing values were found.
        \item The same five sensor features were used for both algorithms.
        \item Logistic Regression was tested with and without feature scaling.
        \item Scaling was done inside a scikit-learn pipeline to avoid data leakage.
    \end{itemize}
\end{frame}

\begin{frame}{Modeling Approach}
    \begin{columns}
        \column{0.48\textwidth}
        \textbf{Logistic Regression}
        \begin{itemize}
            \item Simple baseline model.
            \item Fast to train.
            \item Easy to interpret.
            \item Works well when classes are close to linearly separable.
        \end{itemize}

        \column{0.48\textwidth}
        \textbf{Random Forest}
        \begin{itemize}
            \item Ensemble of decision trees.
            \item Can model nonlinear patterns.
            \item Provides feature importance.
            \item More complex than Logistic Regression.
        \end{itemize}
    \end{columns}
\end{frame}

\begin{frame}{Evaluation Metrics}
    \begin{itemize}
        \item Accuracy: overall proportion of correct predictions.
        \item Precision: how many predicted occupied cases were actually occupied.
        \item Recall: how many actual occupied cases were correctly detected.
        \item F1 score: balance between precision and recall.
        \item Confusion matrix: shows false positives and false negatives directly.
    \end{itemize}
\end{frame}

\begin{frame}{Exploratory Data Analysis}
    \begin{figure}
        \centering
        \includegraphics[width=0.72\textwidth]{correlation_heatmap.png}
        \caption{Feature correlation heatmap.}
    \end{figure}
\end{frame}

\begin{frame}{EDA Takeaways}
    \begin{itemize}
        \item Light has the strongest relationship with occupancy.
        \item CO$_2$ also has a meaningful relationship with occupancy.
        \item Feature scales differ significantly.
        \item Example: Light ranges from 0 to 1546.333, while HumidityRatio stays near 0.004.
        \item This supports testing feature scaling for Logistic Regression.
    \end{itemize}
\end{frame}

\begin{frame}{Validation Results}
    \begin{table}
        \centering
        \caption{Validation Metrics}
        \begin{tabular}{lrrrr}
            \toprule
            Model & Accuracy & Precision & Recall & F1 \\
            \midrule
            Logistic Regression scaled & 0.9771 & 0.9470 & 0.9928 & 0.9694 \\
            Logistic Regression unscaled & 0.9764 & 0.9469 & 0.9907 & 0.9683 \\
            Random Forest & 0.9501 & 0.9458 & 0.9156 & 0.9305 \\
            \bottomrule
        \end{tabular}
    \end{table}

    \vspace{0.2cm}
    Scaled Logistic Regression had the best validation F1 score.
\end{frame}

\begin{frame}{Validation Confusion Matrices}
    \begin{figure}
        \centering
        \includegraphics[width=0.47\textwidth]{validation_confusion_logistic_regression.png}
        \hfill
        \includegraphics[width=0.47\textwidth]{validation_confusion_random_forest.png}
        \caption{Validation confusion matrices: Logistic Regression and Random Forest.}
    \end{figure}
\end{frame}

\begin{frame}{Validation Interpretation}
    \begin{itemize}
        \item Logistic Regression had much higher recall than Random Forest.
        \item Logistic Regression validation confusion matrix:
        \begin{itemize}
            \item 54 false positives
            \item 7 false negatives
        \end{itemize}
        \item Random Forest validation confusion matrix:
        \begin{itemize}
            \item 51 false positives
            \item 82 false negatives
        \end{itemize}
        \item Random Forest missed more occupied cases on the validation split.
    \end{itemize}
\end{frame}

\begin{frame}{Final Holdout Test Results}
    \begin{table}
        \centering
        \caption{Final Holdout Metrics}
        \begin{tabular}{lrrrr}
            \toprule
            Model & Accuracy & Precision & Recall & F1 \\
            \midrule
            Logistic Regression scaled & 0.9925 & 0.9718 & 0.9932 & 0.9824 \\
            Random Forest & 0.9867 & 0.9436 & 0.9961 & 0.9691 \\
            \bottomrule
        \end{tabular}
    \end{table}

    \vspace{0.2cm}
    The final test results confirmed that scaled Logistic Regression was the strongest overall model.
\end{frame}

\begin{frame}{Holdout Confusion Matrices}
    \begin{figure}
        \centering
        \includegraphics[width=0.47\textwidth]{test_confusion_logistic_regression.png}
        \hfill
        \includegraphics[width=0.47\textwidth]{test_confusion_random_forest.png}
        \caption{Holdout confusion matrices: Logistic Regression and Random Forest.}
    \end{figure}
\end{frame}

\begin{frame}{Holdout Interpretation}
    \begin{itemize}
        \item Logistic Regression made fewer total errors on the holdout test set.
        \item Logistic Regression:
        \begin{itemize}
            \item 59 false positives
            \item 14 false negatives
        \end{itemize}
        \item Random Forest:
        \begin{itemize}
            \item 122 false positives
            \item 8 false negatives
        \end{itemize}
        \item Random Forest had slightly higher recall, but lower precision and lower F1.
    \end{itemize}
\end{frame}

\begin{frame}{Random Forest Feature Importance}
    \begin{figure}
        \centering
        \includegraphics[width=0.78\textwidth]{random_forest_feature_importance.png}
        \caption{Random Forest feature importance.}
    \end{figure}
\end{frame}

\begin{frame}{Most Important Features}
    \begin{table}
        \centering
        \caption{Random Forest Feature Importance}
        \begin{tabular}{lr}
            \toprule
            Feature & Importance \\
            \midrule
            Light & 0.6127 \\
            CO$_2$ & 0.2272 \\
            Temperature & 0.0962 \\
            HumidityRatio & 0.0404 \\
            Humidity & 0.0235 \\
            \bottomrule
        \end{tabular}
    \end{table}

    \vspace{0.2cm}
    Light and CO$_2$ were the strongest predictors of occupancy.
\end{frame}

\begin{frame}{Main Findings}
    \begin{itemize}
        \item Scaled Logistic Regression was the best overall model.
        \item The final holdout F1 score was 0.9824.
        \item Random Forest performed well, but did not outperform Logistic Regression.
        \item Light and CO$_2$ were the most important sensor signals.
        \item A simple model was enough for this dataset.
    \end{itemize}
\end{frame}

\begin{frame}{Limitations}
    \begin{itemize}
        \item Only two algorithms were compared.
        \item The timestamp was not engineered into time-based features.
        \item The results come from one dataset and one room environment.
        \item More testing would be needed before deploying the model in other buildings.
    \end{itemize}
\end{frame}

\begin{frame}{Future Work}
    \begin{itemize}
        \item Add time-based features such as hour of day.
        \item Compare more algorithms, such as SVM or Gradient Boosting.
        \item Tune Random Forest hyperparameters more deeply.
        \item Test the model on other rooms or buildings.
        \item Explore real-time IoT deployment.
    \end{itemize}
\end{frame}

\begin{frame}{Conclusion}
    \begin{itemize}
        \item Room occupancy can be predicted accurately using environmental sensor data.
        \item The best model was scaled Logistic Regression.
        \item Final holdout F1 score: \textbf{0.9824}.
        \item Light and CO$_2$ were the strongest indicators of occupancy.
        \item This supports non-invasive occupancy detection for smart-building applications.
    \end{itemize}
\end{frame}

\begin{frame}{References}
    \small
    \begin{enumerate}
        \item Candanedo, L. M., \& Feldheim, V. (2016). Accurate occupancy detection of an office room from light, temperature, humidity and CO$_2$ measurements using statistical learning models. \textit{Energy and Buildings, 112}, 28--39.
        \item UCI Machine Learning Repository. Occupancy Detection Data Set.
        \item Zhang, W., Han, K., Costa, G., \& Li, J. (2022). A review on occupancy prediction through machine learning for enhancing energy efficiency, air quality and thermal comfort in the built environment. \textit{Renewable and Sustainable Energy Reviews}.
    \end{enumerate}
\end{frame}

\begin{frame}
    \centering
    \Huge Questions?
\end{frame}

\end{document}
